\section{Introduction}
\label{sec:intro}

Large language models have turned free text into a queryable attribute.
Systems with semantic operators let an analyst filter, rank, and aggregate records by natural-language criteria that an LLM evaluates tuple by tuple~\cite{lotus,docetl,uqe,thalamusdb}.
Many of the questions asked this way are \emph{decisions over groups}.
A product manager asks which product lines have more than 20\% of reviews reporting a defect (a HAVING query).
A support lead ranks stores by LLM-judged service quality (ORDER BY).
A dashboard shows the five categories with the most negative reviews (Top-$k$).
Evaluating such a query exactly issues one LLM call per tuple.
For the 6.8 million reviews in our evaluation, that is 6.8 million calls and about 860 million input tokens with our prompts, for a single attribute and before the first answer is shown.

Sampling is the natural remedy, and decision queries are where a sampled answer most needs a guarantee.
An estimate that is off by a percentage point is harmless, but a group placed on the wrong side of a threshold or a swapped pair in a ranking is a wrong answer, and nothing in the sample reveals it.
The common practice of sampling until normal-approximation intervals separate provides no such guarantee: the intervals are valid at one pre-planned sample size, and checking them after every batch, as analysts and agents do, inflates the error rate by an amount that depends on the data~\cite{johari2022,ramdas2023savi}.
Certified alternatives exist, but they are expensive.
Range-based bounds with a union bound over checks, as in IFOCUS~\cite{ifocus}, pay for the worst case.
Evergreen uses confidence sequences to stop as soon as a threshold claim is confirmed or refuted, including under sampling without replacement~\cite{evergreen,wsr2024betting}.
A confidence sequence, however, bounds each group's mean on its own: comparing two groups requires both intervals to clear the gap, and Evergreen's estimates do not use the proxy scores that LLM-powered systems already compute.

Proxies are the standard lever for expensive predicates.
BlazeIt and ABae use them to reduce the variance of aggregate estimates~\cite{blazeit,abae}, SUPG to select records with precision or recall targets~\cite{supg}, and LOTUS and BARGAIN to route tuples between cheap and expensive models~\cite{lotus,bargain}.
These methods certify estimates, selections, or tuple-level outputs, mostly at a fixed sample size.
A proxy cannot replace the LLM, however: in our evaluation, answers computed from proxy scores alone were wrong for \numNPROXYWRONG{} of \numNPROXYQ{} queries.
What remains open is how few LLM calls suffice to certify a group-level \emph{decision} when the answer is monitored continuously and a proxy of unknown quality is available.

We answer this question with \sys{}, a decision layer for HAVING, ORDER BY, and Top-$k$ queries over LLM-evaluated attributes.
\sys{} rests on three observations.
First, a decision depends on a few boundary hypotheses, such as ``the mean exceeds the threshold'' or ``group $i$ beats group $j$'', which let \sys{} adapt each bet to the gap relevant to that comparison.
Second, each boundary hypothesis can be tested by betting: the bettor's wealth is a nonnegative supermartingale under the null, and Ville's inequality~\cite{ville1939} bounds the error at every stopping time.
Third, the bets can be placed on proxy-rectified values $y-\beta(p-\bar p_g)$, which have exactly the group mean of the LLM judgments for any proxy, and far smaller variance when the proxy is informative.
\sys{} realizes these observations for sampling without replacement, pairs groups of very different sizes with an exactly unbiased per-draw statistic, and stops each group as soon as its decisions are certified.

This paper makes the following contributions.
\begin{itemize}
\item We define certified decision queries over LLM-evaluated attributes, with an indifference margin that makes certification possible before every tuple is judged (\Cref{sec:problem}).
\item We design proxy-rectified boundary e-processes for single groups and for pairs under sampling without replacement, and prove that \sys{} returns an incorrect answer with probability at most $\alpha$ under any stopping rule and any proxy (\Cref{sec:method}).
\item We compile HAVING, ORDER BY, and Top-$k$ queries into these tests with per-group adaptive stopping (\Cref{sec:processing}).
\item We evaluate \sys{} on \numNQ{} queries over three datasets, two of them judged by a real LLM. \sys{} never returned an incorrect answer, used \numRWSR$\times$ fewer LLM calls than Evergreen-style confidence sequences (up to \numRWSRMAX$\times$) and \numRHS--\numRBSP$\times$ fewer than IFOCUS- and BlazeIt-style bounds, and needed about as many calls as uncertified sampling, which erred in up to 13\% of executions. With two LLMs in the loop, latency followed the number of calls, and HAVING queries ran \numLiveHavingMin$\times$ and \numLiveHavingMax$\times$ faster than full evaluation (\Cref{sec:exp}).
\end{itemize}
